% !TEX root = ../matan.tex

\section{Функциональные последовательности и ряды. Поточечная и равномерная сходимости семейства функций. Определения, примеры.}

\begin{Definition}{Последовательность и ряд функций}{}
	Пусть $\Omega$~--- множество и $f_k \colon \Omega \to \C$. Тогда $\{f_k\}_{k \in \NN}$~--- \textbf{последовательность функций}; для каждой фиксированной точки $x \in \Omega$ возникает числовая последовательность $\{f_k(x)\}$.

	Формальное выражение $\sum_{k=1}^{\infty} f_k$ называется \textbf{функциональным рядом}; его частичные суммы $S_n(x) = \sum_{k=1}^{n} f_k(x)$.
\end{Definition}

\begin{Definition}{Поточечная сходимость}{}
	Последовательность $f_k \colon \Omega \to \C$ \textbf{сходится поточечно} к $f \colon \Omega \to \C$, если для каждого $x \in \Omega$ числовая последовательность $f_k(x)$ сходится к $f(x)$:
	\[
	\forall x \in \Omega \colon \quad \lim_{k \to \infty} f_k(x) = f(x).
	\]
	Ряд $\sum_{k=1}^{\infty} f_k$ сходится поточечно к $S$, если $S_n(x) \to S(x)$ для каждого $x \in \Omega$.
\end{Definition}

\begin{Remark}{}{}
	В поточечной сходимости номер $N$, начиная с которого $|f_k(x) - f(x)| < \eps$, зависит и от $\eps$, и от точки $x$: $N = N(\eps, x)$.
\end{Remark}

\begin{Definition}[required]{Равномерная сходимость}{}
	Последовательность $f_n \colon \Omega \to \C$ \textbf{сходится равномерно} к $f \colon \Omega \to \C$, если
	\[
	\forall \eps > 0\ \exists N\ \forall n > N\ \forall x \in \Omega \colon \quad |f_n(x) - f(x)| < \eps.
	\]
	Эквивалентно, в терминах равномерной нормы:
	\[
	\|f_n - f\|_{\infty} = \sup_{x \in \Omega} |f_n(x) - f(x)| \longrightarrow 0 \qquad (n \to \infty).
	\]
	Ряд $\sum_{k=1}^{\infty} f_k$ сходится равномерно, если равномерно сходится последовательность его частичных сумм $S_n$.
\end{Definition}

Ключевое отличие от поточечной сходимости: здесь номер $N = N(\eps)$ годится сразу для всех $x \in \Omega$.

\begin{Statement}{Равномерная сходимость влечёт поточечную}{}
	Если $f_n \to f$ равномерно на $\Omega$, то $f_n \to f$ поточечно.
\end{Statement}

\begin{proof}
	По равномерной сходимости для любого $\eps > 0$ существует $N$ такое, что $\sup_{x \in \Omega} |f_n(x) - f(x)| < \eps$ при $n > N$. Тогда тем более для каждого фиксированного $x$ имеем $|f_n(x) - f(x)| < \eps$ при $n > N$, то есть $f_n(x) \to f(x)$.
\end{proof}

\begin{Remark}{Обратное неверно}{}
	Поточечная сходимость не влечёт равномерную. Например, $f_n(x) = x^n$ на $(0, 1)$ сходится поточечно к $0$, но $\sup_{x \in (0,1)} |x^n - 0| = 1$ при каждом $n$, поэтому сходимость не равномерна.
\end{Remark}

\begin{Example}{Примеры}{}
	\begin{enumerate}
		\item $f_n(x) = \dfrac{x}{n}$ на $[0, 1]$. Поточечно $f_n \to 0$; при этом $\sup_{[0,1]} |x/n| = 1/n \to 0$, значит сходимость \emph{равномерна}.
		\item $f_n(x) = x^n$ на $[0, 1]$. Поточечный предел
		\[
		f(x) = \begin{cases} 0, & 0 \le x < 1, \\ 1, & x = 1, \end{cases}
		\]
		разрывен, тогда как все $f_n$ непрерывны; сходимость \emph{не равномерна} (предел равномерно сходящейся последовательности непрерывных функций непрерывен).
		\item Ряд $\sum_{k=0}^{\infty} x^k$ на $[0, r]$ при $0 < r < 1$ сходится равномерно к $\dfrac{1}{1 - x}$ (мажорируется сходящимся числовым рядом $\sum r^k$), но на всём $(-1, 1)$~--- лишь поточечно.
	\end{enumerate}
\end{Example}
